%!TEX program = uplatex
\RequirePackage{plautopatch}
\documentclass[uplatex,dvipdfmx]{jlreq}
\usepackage{amsmath,amssymb,amsfonts,amsthm}
\usepackage{enumerate}

\pagestyle{empty}

\begin{document}

%======================================================================
% 講演1：加藤毅「分子生物学の非常に基礎的な事柄についての解説」
%======================================================================
\begin{center}
{\large\bfseries 分子生物学の非常に基礎的な事柄についての解説}\\[6pt]
\hspace*{\fill}{\large\bfseries 加藤 毅（京都大学 大学院理学研究科・数学教室）}
\end{center}

今回の
ENCOUNTER with MATHEMATICS
はけいはんなでこれまで行われてきた
「数学者のための分子生物学入門」のダイジェスト版といえる。
ここでは分子生物学の専門家によるそれぞれの講義を
理解するための基礎事項を素人の立場から解説する。

\begin{center}\textsc{参考文献}\end{center}

上野健爾・加藤毅編、数学者のための分子生物学入門
1（2003）、2（2004）、3（2005）、
4（2006）作成中

入手方法：
京都大学理学部数学教室事務　田中さん
(mugen@math.kyoto-u.ac.jp)
または
加藤
(tkato@math.kyoto-u.ac.jp)
まで連絡のこと。

1、3については、
雑誌「物性研究」2003-10 81-1, 2005-10 85-1
に同じものが載っています。

\vfill

\newpage

%======================================================================
% 講演2：阿久津達也「生体内ネットワークの数理モデル」
%======================================================================
\begin{center}
{\large\bfseries 生体内ネットワークの数理モデル}\\[6pt]
\hspace*{\fill}{\large\bfseries 阿久津 達也（京都大学・化学研究所・バイオインフォマティクスセンター）}
\end{center}

細胞内で遺伝子、タンパク質、化合物が複雑に相互作用しあ
うことにより生命活動が維持されていますが、それらの関係は、
いくつかの種類のネットワーク構造（グラフ構造）としてモデ
ル化されています。そして、データベース解析により、これら
のネットワーク構造の多くが、スケールフリー性（頂点の次数
分布のべき乗則）を持つということが示唆されてきました。

また、生体内ネットワークのみならず、インターネットや論
文の共著関係のなすネットワークなど、人工的、社会的なネッ
トワークの多くもスケールフリー性を持つことが示唆されてき
ました。一方、スケールフリー性やその他の特徴を持つネット
ワークを人工的に生成するために、様々な数理モデルが提案さ
れてきました。

今回は、スケールフリー性を中心にして、生体内ネットワー
クに関する数理モデルの概要と私たちの研究についてお話しし
ようと思います。

\begin{center}\textsc{参考文献}\end{center}
\begin{enumerate}[{[1]}]
  \item 増田直紀、今野紀雄：複雑ネットワークの科学、産業図書、2005.
  \item 阿久津達也, 落合友四郎, Jose~C.~Nacher：
    生物情報ネットワークの構造およびダイナミクス解析、
    蛋白質・核酸・酵素、Vol.~50, No.~16, pp.~2288--2293, 2005.
\end{enumerate}

\vfill

\newpage

%======================================================================
% 講演3：岡本祐幸「タンパク質の折り畳みに関する最適化問題」
%======================================================================
\begin{center}
{\large\bfseries タンパク質の折り畳みに関する最適化問題}\\[6pt]
\hspace*{\fill}{\large\bfseries 岡本 祐幸（名古屋大学 大学院理学研究科・物理学教室）}
\end{center}

タンパク質は自然の環境下で、一定の立体構造に折り畳みます。
これは、数学的には、ある評価関数（系のエネルギー関数）の
最適化問題であるということができます。
すなわち、多変数関数の最小値を求める問題です。
ここでは、タンパク質の自然の立体構造をランダムな
初期構造からスタートする計算機シミュレーションに
よって求める問題について、
幾つかの有効な手法を紹介します。

\begin{center}\textsc{参考文献}\end{center}

岡崎進・岡本祐幸（編）「化学フロンティア」No.~8（化学同人, 2002）第2章

\vfill

\newpage

%======================================================================
% 講演4：斎藤成也「遺伝子の系統樹と系統ネットワーク」
%======================================================================
\begin{center}
{\large\bfseries 遺伝子の系統樹と系統ネットワーク}\\[6pt]
\hspace*{\fill}{\large\bfseries 斎藤 成也（国立遺伝学研究所）}
\end{center}

遺伝子の系統樹は、
DNAの半保存的複製の繰り返しによって生じるので、
基本は二分岐の有根系統樹である。
しかし、DNAはヌクレオチドが長くつながった分子なので、
ある部分と他の部分の系統関係が異なってしまう場合が生じる。
これは、組換え（recombination）、
遺伝子変換（gene conversion）、
ドメインの融合・分離などの現象によるものだ。

そこで、塩基配列間の相互関係を系統樹であると
頭から決めてかからず、配列データに即して記述する必要がある。
これが「系統ネットワーク」である。
系統ネットワークは、すべての節（node）が結合しているグラフである。
グラフ理論では、節間の結合を表すのに線 (line, edge) という
用語を用いるが、これには長さの概念がない。
ところが系統ネットワークでは、
2個の節（塩基配列やアミノ酸配列など）の間の
長さ（蓄積した突然変異の個数）も考慮する。

本講義では、
雑誌「数理科学」に以前掲載した内容（斎藤成也、2004--2005）を
もとにして、系統樹と系統ネットワークについて論じる。

\begin{center}\textsc{参考文献}\end{center}
\begin{enumerate}[{[1]}]
  \item 斎藤成也 (2004) 系統樹とは（その1）.
    連載「ゲノム進化学の展開」第2回.
    数理科学、42巻9号、68--74頁.
  \item 斎藤成也 (2004) 系統樹とは（その2）.
    連載「ゲノム進化学の展開」第3回.
    数理科学、42巻10号、72--78頁.
  \item 斎藤成也 (2005) 形質状態データから系統樹を作成する方法.
    連載「ゲノム進化学の展開」第5回.
    数理科学、43巻1号、64--72頁.
  \item 斎藤成也 (2005) 系統ネットワークを用いた配列間関係の表現.
    連載「ゲノム進化学の展開」第6回.
    数理科学、43巻3号、60--66頁.
\end{enumerate}

\vfill

\newpage

%======================================================================
% 講演5：田中博「蛋白質間相互作用ネットワークの構造と進化」
%======================================================================
\begin{center}
{\large\bfseries 蛋白質間相互作用ネットワークの構造と進化}\\[6pt]
\hspace*{\fill}{\large\bfseries 田中 博（東京医科歯科大学 大学院疾患生命科学研究部）}
\end{center}

たんぱく質間相互作用の大規模スクリーニングの進展に伴い、
生命現象をネットワークの視点で捉える試みがなされている。
その結果、たんぱく質間相互作用ネットワーク中には、
生命系に特徴的な構造が存在しており（階層構造、クリーク構造）、
それらは特定の機能を担うだけでなく進化的にも良く保存されている
ことが判明している。
しかしながら、
従来の遺伝子重複のみを基にしたネットワーク進化モデルでは、
例えば、クリーク構造の存在を充分に説明することは出来ない。

我々の進化モデルでは遺伝子重複をベースとしつつも、
2ステップ離れた関係にあるタンパク質間に
新しい相互作用を加える過程を導入することで、
クリーク構造の再現に成功した。

\vfill

\end{document}
